% A submission that cites made-up works and a retracted paper.
\documentclass[conference]{IEEEtran}
\begin{document}
\title{FlakeScope: Ranking Flaky Test Failures}
\author{\IEEEauthorblockN{Anonymous Author(s)}}
\maketitle
\begin{abstract}
Flaky tests pass and fail on the same code. FlakeScope ranks the failures of a build by the chance that each
failure is flaky and places the flaky failure first in 71\% of 4,210 failures of ten projects.
\end{abstract}
\section{Introduction}
Flaky tests pass and fail on the same code. Such tests waste the time of developers, hide real faults, and erode the
trust in continuous integration. This paper presents FlakeScope, a tool that ranks the tests of a project by the
chance that a failure is flaky. FlakeScope reads the history of the test runs and the changes of each commit.

Our evaluation covers ten open source projects with 4,210 test failures. FlakeScope places the flaky failure first
in 71\% of the cases, compared with 38\% for a rerun of the failed tests alone.

\section{Background}
A test is flaky when it passes and fails on the same version of the code. Common causes are asynchronous waits,
shared state between tests, and the order in which the tests run. A rerun of each failure detects some flaky
tests. However, the rerun costs time on every build.

\section{Approach}
FlakeScope collects three signals for each failure. The first signal is the history of the test, such as the
fraction of past failures that a rerun fixed. The second signal is the distance between the failing test and the changed
code. The third signal is the order of the tests in the failing run. A gradient boosted model combines the three
signals into one score.

\section{Evaluation}
We collected the build logs of ten projects over two years. Two annotators labeled each failure as flaky or real
and agreed on 92\% of the failures before a discussion of the rest. FlakeScope ranks 2,988 of 4,210 failures
correctly at the first position. The history signal contributes the most, followed by the distance signal.

\subsection{Projects}
The ten projects span build tools, web frameworks, databases, and libraries for data processing. Each project
runs its tests on every commit in a public continuous integration service and keeps the logs of at least two
years. The projects hold between 1,200 and 18,000 tests and between 40 and 900 test failures per year. Seven
projects run their tests in parallel, which makes failures that depend on the order of the tests more common.

\subsection{Accuracy of the Ranking}
For each failed build, we sort the failures by the score of FlakeScope and record the position of the first
flaky failure. FlakeScope places the flaky failure first in 71\% of the builds and within the first three
positions in 89\% of the builds. A rerun of the failed tests places the flaky failure first in 38\% of the
builds, since a rerun misses flaky tests that fail again by chance. The accuracy stays above 65\% in every
project, and the project with the most parallel tests reaches the highest accuracy.

\subsection{Contribution of the Signals}
We remove each signal in turn and repeat the ranking. Without the history signal, the accuracy drops to 52\%.
Without the distance signal, it drops to 63\%. Without the order signal, it drops to 68\%. The three signals
therefore complement each other, and the history of a test carries the most information about its flakiness.

\subsection{Cost}
FlakeScope scores the failures of a build in less than two seconds on a laptop. The model trains in four minutes
on the history of one project. A rerun of every failed test, in contrast, adds a median of eleven minutes to a
failed build in the ten projects.

\section{Discussion}
Developers can use the ranking in two ways. First, the ranking can order the failures in the report of a build,
such that the failures that need attention come first. Second, a project can rerun only the failures that rank
low, which saves the time of the reruns that a flaky failure would waste. Both uses keep a developer in the loop,
since a ranking never hides a failure.

The ranking also helps the maintainers of a project to find the tests that fail most often for reasons outside
the code under test. These tests are candidates for a fix of their waits, their shared state, or their order.

\section{Threats to Validity}
The labels of the failures come from two annotators who read the logs and the code. Their agreement of 92\%
indicates a reliable labeling, but some failures remain hard to judge from a log alone. The ten projects are open
source projects with public logs, and industrial projects may show other causes of flakiness. The model of
FlakeScope trains on the history of each project, therefore a new project needs a few weeks of history before
the ranking becomes accurate.

\section{Related Work}
Prior work studies the causes of flaky tests and detects order-dependent tests. FlakeScope complements these
detectors with a ranking that a developer can read after each failed build.

\section{Conclusion}
FlakeScope ranks the failures of a build by the chance that each failure is flaky. A tool paper with the full
evaluation and the data set follows.

\begin{thebibliography}{3}
\bibitem{fake1} M. Quill and R. Stone, ``Quantum scheduling of flaky test reruns at hyperscale,'' in
\emph{Proc. ICSE}, 2023, pp. 1--12.
\bibitem{fake2} T. Arden, ``FlakyNet: a foundation model for every failed build,'' \emph{IEEE Trans. Softw. Eng.},
vol. 51, no. 4, 2025.
\bibitem{retracted} A. J. Wakefield \emph{et al.}, ``Ileal-lymphoid-nodular hyperplasia, non-specific colitis, and
pervasive developmental disorder in children,'' \emph{The Lancet}, vol. 351, no. 9103, pp. 637--641, 1998,
doi:10.1016/S0140-6736(97)11096-0.
\end{thebibliography}
\end{document}
